% =====================================================================
% Rapport du projet « API REST de produits »
% Compilation :  cd rapport && tectonic rapport.tex     (ou : xelatex rapport.tex, deux fois)
% Les captures d'écran sont facultatives : déposer les images dans figures/
% (noms attendus dans docs/GUIDE-PRESENTATION.md) ; elles apparaissent automatiquement.
% =====================================================================
\documentclass[11pt,a4paper]{article}

\usepackage[T1]{fontenc}
\usepackage[french]{babel}
\usepackage[margin=2.5cm]{geometry}
\usepackage{lmodern}
\usepackage{microtype}
\usepackage{graphicx}
\usepackage{float}
\usepackage{booktabs}
\usepackage{tabularx}
\usepackage{array}
\usepackage[export]{adjustbox}
\usepackage[table]{xcolor}
\usepackage{multirow}
\usepackage{enumitem}
\usepackage{fancyhdr}
\usepackage{fancyvrb}
\usepackage{upquote}
\usepackage{caption}
\usepackage{tikz}
\usetikzlibrary{positioning,arrows.meta,fit,backgrounds,shapes.misc}
\usepackage[hidelinks]{hyperref}

\setlength{\emergencystretch}{3em}
% sommaire compact : tient sur une seule page
\usepackage{tocloft}
\setlength{\cftbeforesecskip}{0.45em}
\setlength{\cftbeforesubsecskip}{0pt}
\setlist{itemsep=2pt,topsep=4pt}
\renewcommand{\arraystretch}{1.2}
\captionsetup{font=small,labelfont=bf}

% ---------- Couleurs et macros
\definecolor{bleu}{HTML}{404040}
\definecolor{gris}{HTML}{F2F2F2}
\definecolor{vert}{HTML}{000000}
\definecolor{rouge}{HTML}{595959}
\newcommand{\code}[1]{\texttt{\detokenize{#1}}}
\newcommand{\aremplir}[1]{\textcolor{rouge}{\textbf{[#1]}}}
\newcolumntype{L}[1]{>{\raggedright\arraybackslash}p{#1}}
\newcolumntype{Y}{>{\raggedright\arraybackslash}X}
\DefineVerbatimEnvironment{Code}{Verbatim}{fontsize=\footnotesize,frame=single,framesep=2mm,rulecolor=\color{gray!60},formatcom=\color{black}}

% Capture d'écran facultative : n'affiche rien si le fichier est absent
\newcommand{\capture}[3][0.9]{%
  \IfFileExists{figures/#2}{%
    \begin{figure}[H]\centering
      \includegraphics[width=#1\textwidth]{figures/#2}
      \caption{#3}
    \end{figure}}{}}

% ---------- En-têtes
\pagestyle{fancy}
\fancyhf{}
\renewcommand{\headrulewidth}{0.3pt}
\fancyhead[L]{\small Développement à base de composants et services web}
\fancyhead[R]{\small Groupe~1 --- M1 IBAM 2026}
\fancyfoot[C]{\small\thepage}

% ---------- Styles TikZ
\tikzset{
  boite/.style={draw=bleu,fill=bleu!8,rounded corners=2pt,minimum height=1.1cm,
                text width=2.4cm,align=center,font=\footnotesize,inner sep=3pt},
  ext/.style={boite,draw=gray,fill=gris},
  flux/.style={-{Stealth[length=2.2mm]},thick,color=bleu!80!black},
  etiquette/.style={font=\scriptsize\itshape,color=black!70,midway,fill=white,inner sep=1pt}
}

\begin{document}

% =====================================================================
% PAGE DE GARDE
% =====================================================================
\begin{titlepage}
\centering
\vspace*{1cm}
{\large Master 1 --- IBAM \hfill Année universitaire 2025--2026\par}
\vspace{0.4cm}
{\large Développement à base de composants et services web\par}
\vspace{2.2cm}
\rule{\textwidth}{0.6pt}\\[0.5cm]
{\Huge\bfseries API REST de produits\par}
\vspace{0.4cm}
{\Large Bases de données, ETL Apache Hop,\\ API Spring Boot et supervision Prometheus\par}
\rule{\textwidth}{0.6pt}\\[1.2cm]
{\large\bfseries Rapport de projet --- Exercice 1\par}
\vspace{2cm}
\begin{tabular}{rl}
  \textbf{Groupe} & 1 \\[4pt]
  \textbf{Membres} & BOUGOUMA Leïla --- partie 1 : bases de données \\
                   & OUARMA Landry --- partie 2 : ETL \\
                   & WAONGO Osiris --- partie 3 : API \\
                   & ZAPRE Malika --- partie 4 : supervision \\
                   & KOURA Stéphane --- partie 5 : tests \\[4pt]
  \textbf{Enseignant} & Mr Oumar KY \\[4pt]
  \textbf{Date} & \today
\end{tabular}
\vfill
{\small Code source, scripts et configuration : dépôt Git du projet (branche \code{main}).\par}
\end{titlepage}

\begin{abstract}
\noindent
Ce rapport présente la réalisation d'une chaîne complète de traitement de données de produits.
Des données brutes et volontairement imparfaites, stockées dans MySQL, sont contrôlées, normalisées
et dédoublonnées par un ETL Apache Hop, puis chargées dans PostgreSQL. Une API REST Spring Boot expose
ces données (consultation paginée, création, modification, suppression) avec un contrat OpenAPI,
une validation stricte et un format d'erreur unique. Prometheus supervise l'API en continu.
L'ensemble est décrit par un seul fichier Docker Compose et se lance sur n'importe quelle machine
disposant de Docker. Sur le jeu de test de 40 lignes, l'ETL charge 29 produits et en rejette 10 avec
un motif précis ; l'API est couverte par 27 tests automatisés et par une collection Bruno de 18 requêtes.
\end{abstract}
\clearpage

\tableofcontents
\clearpage

% =====================================================================
\section{Introduction}
% =====================================================================

\subsection{Contexte et objectif}

Dans une organisation, les données de produits proviennent souvent de sources hétérogènes et peu fiables :
champs vides, prix négatifs, catégories écrites de dix façons différentes, doublons. Les exposer telles quelles
à des applications clientes les rend inutilisables. Ce projet, réalisé dans le cadre du module
« Développement à base de composants et services web », construit la chaîne qui permet de passer de ces
données brutes à un service web fiable et supervisé.

La chaîne comporte cinq parties, confiées chacune à un membre du groupe :
\begin{enumerate}
  \item les \textbf{bases de données} : une base source MySQL contenant des données sales, une base cible
        PostgreSQL contenant des données propres, et un jeu de test de 40 lignes ;
  \item l'\textbf{ETL} (extraction, transformation, chargement) qui alimente PostgreSQL à partir de MySQL ;
  \item l'\textbf{API REST} Spring Boot qui expose les produits ;
  \item la \textbf{supervision} de l'API avec Prometheus ;
  \item les \textbf{tests} et la validation de bout en bout.
\end{enumerate}

\subsection{Principes directeurs}

Trois principes ont guidé toutes les décisions techniques.
\begin{itemize}
  \item \textbf{Reproductibilité} : tout se lance avec Docker Compose ; aucune installation manuelle
        d'une base de données ou d'un ETL n'est nécessaire.
  \item \textbf{Aucun secret dans le dépôt} : les mots de passe sont lus dans un fichier \code{.env}
        non versionné ; seul un modèle \code{.env.example} est partagé.
  \item \textbf{Contrôle de la qualité à chaque étape} : l'ETL rejette les données invalides, la base cible
        les refuse par des contraintes, et l'API les refuse par une validation explicite.
\end{itemize}

% =====================================================================
\section{Architecture générale}
% =====================================================================

\subsection{Vue d'ensemble}

La figure~\ref{fig:archi} montre les composants et les flux de données. Les données circulent de la gauche
vers la droite : MySQL (source), ETL, PostgreSQL (cible), puis API. Prometheus interroge l'API à intervalle
régulier. Tous les composants tournent dans des conteneurs Docker décrits par un unique fichier
\code{docker-compose.yml}.

\begin{figure}[H]
\centering
\begin{adjustbox}{max width=\textwidth}
\begin{tikzpicture}[node distance=1.4cm]
  \node[ext]  (mysql) {MySQL 8.0\\ base \emph{source}\\ données brutes};
  \node[boite,right=of mysql] (etl) {ETL\\ Apache Hop 2.18.1};
  \node[ext,right=of etl]     (pg)  {PostgreSQL 16\\ base \emph{cible}\\ données propres};
  \node[boite,right=of pg]    (api) {API REST\\ Spring Boot 3.3\\ Java 17};
  \node[ext,above=1.5cm of api]  (cli)  {Clients\\ Swagger UI, Bruno,\\ navigateur};
  \node[boite,below=1.5cm of api] (prom) {Prometheus 2.53\\ supervision};

  \draw[flux] (mysql) -- node[etiquette,above=2pt]{lecture} (etl);
  \draw[flux] (etl)   -- node[etiquette,above=2pt]{écriture} (pg);
  \draw[flux] (api)   -- node[etiquette,above=2pt]{JPA / JDBC} (pg);
  \draw[flux] (cli)   -- node[etiquette,right=2pt]{HTTP / JSON} (api);
  \draw[flux,dashed] (prom) -- node[etiquette,right=2pt]{collecte toutes les 5\,s} (api);

  \begin{scope}[on background layer]
    \node[draw=gray,dashed,rounded corners=4pt,fill=gray!4,inner sep=10pt,
          fit=(mysql)(etl)(pg)(api)(cli)(prom)] (cadre) {};
    \node[font=\scriptsize,gray,anchor=north west] at ([xshift=6pt,yshift=-5pt]cadre.north west)
          {Docker Compose --- réseau interne};
  \end{scope}
\end{tikzpicture}
\end{adjustbox}
\caption{Architecture générale du projet.}
\label{fig:archi}
\end{figure}

\subsection{Les services Docker}

\begin{table}[H]
\centering\small
\begin{tabularx}{\textwidth}{|L{2.2cm}|Y|L{2.5cm}|Y|}
\hline
\rowcolor{bleu!15}
\textbf{Service} & \textbf{Image} & \textbf{Port hôte} & \textbf{Rôle} \\
\hline
\code{mysql}      & \code{mysql:8.0}            & 3306 & Base source, initialisée par des scripts SQL \\
\hline
\code{postgres}   & \code{postgres:16}          & 5432 (\code{PG_PORT}) & Base cible, utilisée par l'ETL et l'API \\
\hline
\code{etl}        & Apache Hop 2.18.1 (image construite dans \code{etl/}) & aucun & Exécute le workflow puis s'arrête (à la demande) \\
\hline
\code{hop-web}    & idem                        & 8090 & Interface graphique de Hop (optionnelle) \\
\hline
\code{api}        & construite depuis \code{api/} & 8080 & API REST Spring Boot \\
\hline
\code{prometheus} & \code{prom/prometheus:v2.53.0} & 9090 & Collecte et consultation des métriques \\
\hline
\end{tabularx}
\caption{Services définis dans \code{docker-compose.yml}.}
\end{table}

Plusieurs mécanismes rendent le démarrage fiable :
\begin{itemize}
  \item des \textbf{contrôles de santé} (\emph{healthchecks}) sur MySQL et PostgreSQL ; l'API et l'ETL n'attendent
        que les bases soient \code{healthy} avant de démarrer (\code{depends_on} avec condition) ;
  \item des \textbf{volumes nommés} (\code{mysql_data}, \code{pg_data}) qui conservent les données entre deux
        arrêts ; \code{docker compose down -v} réinitialise tout ;
  \item des \textbf{scripts d'initialisation} montés dans \code{/docker-entrypoint-initdb.d}, exécutés
        automatiquement au premier démarrage, dans l'ordre alphabétique (\code{01-}, \code{02-}) ;
  \item des \textbf{ports configurables} dans \code{.env}, utiles lorsqu'une machine héberge déjà une base
        de données (voir la section~\ref{sec:difficultes}).
\end{itemize}

\subsection{Choix structurants}

\begin{description}
  \item[Le schéma est écrit en SQL, pas généré par Hibernate.] L'API démarre avec
        \code{ddl-auto: validate} : si la table ne correspond pas au code, elle refuse de démarrer. Le même
        script SQL sert donc de référence à l'ETL, à l'API et aux tests.
  \item[Deux bases distinctes.] La séparation source/cible reproduit un cas réel : la source est en
        lecture seule pour l'ETL et ne contient aucune contrainte, la cible impose toutes les règles.
  \item[Un seul réseau Docker.] Les conteneurs se joignent par leur nom (\code{mysql}, \code{postgres},
        \code{api}) ; seuls les ports utiles sont publiés sur la machine.
\end{description}

% =====================================================================
\section{Partie 1 --- Bases de données et jeu de test}
% =====================================================================

\subsection{Base source MySQL}

La table \code{produits_source} est créée \emph{sans aucune contrainte} (ni \code{NOT NULL}, ni
\code{CHECK}), de façon volontaire : elle représente une source externe peu fiable. C'est à l'ETL de faire
le contrôle qualité.

\begin{table}[H]
\centering\small
\begin{tabular}{|l|l|p{7.5cm}|}
\hline
\rowcolor{bleu!15}
\textbf{Colonne} & \textbf{Type} & \textbf{Remarque} \\
\hline
\code{id}            & \code{INT AUTO_INCREMENT} & Clé primaire \\
\hline
\code{nom}           & \code{VARCHAR(150)}       & Peut être vide ou nul \\
\hline
\code{description}   & \code{TEXT}               & \\
\hline
\code{prix}          & \code{DECIMAL(10,2)}      & Peut être négatif ou nul \\
\hline
\code{quantite}      & \code{INT}                & Peut être négative \\
\hline
\code{categorie}     & \code{VARCHAR(100)}       & Écrite de plusieurs façons (\texttt{electro}, \texttt{ÉLECTRONIQUE}\ldots) \\
\hline
\code{date_creation} & \code{DATETIME}           & \code{DEFAULT CURRENT_TIMESTAMP} \\
\hline
\end{tabular}
\caption{Structure de la base source.}
\end{table}

La base est configurée en \code{utf8mb4} afin de stocker correctement les accents (« Électronique »,
« Vêtements »).

\subsection{Base cible PostgreSQL}

La base cible contient deux tables et applique toutes les règles de gestion.

\begin{table}[H]
\centering\small
\begin{tabularx}{\textwidth}{|L{3.1cm}|L{3.8cm}|Y|}
\hline
\rowcolor{bleu!15}
\textbf{Table} & \textbf{Colonne} & \textbf{Contrainte} \\
\hline
\multirow{7}{*}{\code{produits}}
  & \code{id}            & \code{SERIAL PRIMARY KEY} \\ \cline{2-3}
  & \code{nom}           & \code{VARCHAR(150) NOT NULL} \\ \cline{2-3}
  & \code{description}   & \code{TEXT} \\ \cline{2-3}
  & \code{prix}          & \code{NUMERIC(10,2) NOT NULL CHECK (prix >= 0)} \\ \cline{2-3}
  & \code{quantite}      & \code{INTEGER NOT NULL DEFAULT 0 CHECK (quantite >= 0)} \\ \cline{2-3}
  & \code{categorie}     & \code{VARCHAR(100) NOT NULL} \\ \cline{2-3}
  & \code{date_creation} & \code{TIMESTAMP NOT NULL DEFAULT NOW()} \\
\hline
\multirow{3}{*}{\code{produits_rejets}}
  & \code{id_source}  & \code{INT} (identifiant dans MySQL) \\ \cline{2-3}
  & \code{motif}      & \code{VARCHAR(255)} (raison du rejet) \\ \cline{2-3}
  & \code{date_rejet} & \code{TIMESTAMP DEFAULT NOW()} \\
\hline
\end{tabularx}
\caption{Structure de la base cible.}
\end{table}

Deux index accélèrent les requêtes fréquentes : un index sur \code{categorie} (filtre de l'API) et un index
sur \code{(nom, categorie)}, qui correspond à la clé fonctionnelle utilisée par l'ETL pour détecter les
doublons.

\subsection{Jeu de test}
\label{sec:jeu-ref}

Le jeu de test insère 40 produits, construits pour exercer chaque règle de l'ETL
(tableau~\ref{tab:jeu}). Le fichier \code{db/DONNEES-TEST.md} détaille chaque ligne et sert de référence aux
parties ETL et tests.

\begin{table}[H]
\centering\small
\begin{tabularx}{\textwidth}{|L{2.3cm}|Y|L{4.2cm}|}
\hline
\rowcolor{bleu!15}
\textbf{Lignes} & \textbf{Contenu} & \textbf{Résultat attendu} \\
\hline
1 à 29     & Produits valides, catégories écrites de plusieurs façons & Chargées dans \code{produits} \\
\hline
30         & Doublon de la ligne 7 (\emph{Souris sans fil}, Informatique) & Mise à jour de la ligne existante \\
\hline
31 à 38, 40 & Nom, prix, quantité ou catégorie manquant ou invalide & Rejetées avec un motif \\
\hline
39         & Catégorie inconnue (\texttt{jouets}) & Rejetée (décision du groupe) \\
\hline
\end{tabularx}
\caption{Répartition des 40 lignes de test.}
\label{tab:jeu}
\end{table}

Le jeu contient aussi des cas limites \emph{valides} qu'il ne faut pas rejeter : une quantité égale à zéro
(ligne~21) et un prix égal à zéro (ligne~27).

Le script \code{db/verifier-bases.sh} contrôle que les deux bases répondent et affiche le nombre de lignes
(40 dans la source ; les tables cibles sont vides avant le premier passage de l'ETL).

% =====================================================================
\section{Partie 2 --- ETL avec Apache Hop}
% =====================================================================

\subsection{Choix de l'outil}

L'énoncé proposait Talend Open Studio. Cet outil n'est plus maintenu depuis 2024 et son installation
n'est plus fiable ; l'énoncé autorisant une autre solution ETL, le groupe a retenu \textbf{Apache Hop 2.18.1}
pour trois raisons : c'est un projet libre et actif de la fondation Apache (issu de Pentaho Kettle) ; ses
traitements sont enregistrés en XML et donc versionnables avec Git ; il fournit une image Docker officielle,
ce qui rend l'exécution reproductible.

\subsection{Le pipeline de traitement}

Le pipeline \code{pl_produits} traite chaque ligne en cinq étapes (figure~\ref{fig:etl}).

\begin{figure}[H]
\centering
\begin{adjustbox}{max width=\textwidth}
\begin{tikzpicture}[node distance=0.4cm,boite/.append style={text width=2.35cm,minimum height=1.5cm}]
  \node[ext]  (a) {1.~Extraction\\ \texttt{produits\_\allowbreak source}};
  \node[boite,right=of a] (b) {2.~Normalisation\\ trim, minuscules, sans accents};
  \node[boite,right=of b] (c) {3.~Correspondance\\ référentiel de catégories};
  \node[boite,right=of c] (d) {4.~Contrôle\\ règles de gestion};
  \node[ext,right=of d]   (e) {5.~Chargement\\ \texttt{produits}};
  \node[ext,below=1.1cm of d] (r) {\texttt{produits\_\allowbreak rejets}\\ avec motif};
  \draw[flux] (a)--(b); \draw[flux] (b)--(c); \draw[flux] (c)--(d);
  \draw[flux] (d)--node[etiquette,above=2pt]{valide}(e);
  \draw[flux,color=rouge] (d)--node[etiquette,right=2pt]{invalide}(r);
\end{tikzpicture}
\end{adjustbox}
\caption{Étapes du pipeline ETL.}
\label{fig:etl}
\end{figure}

\begin{enumerate}
  \item \textbf{Extraction} : lecture de la table \code{produits_source} dans MySQL.
  \item \textbf{Normalisation} : suppression des espaces superflus dans le nom ; mise en minuscules de la
        catégorie, suppression des espaces et des accents (« ÉLECTRONIQUE » devient « electronique »).
  \item \textbf{Correspondance} : recherche de la catégorie officielle dans le référentiel
        \code{categories.csv} (« electro » devient « Electronique »). Il y a cinq catégories officielles :
        Electronique, Informatique, Vetements, Alimentation et Maison.
  \item \textbf{Contrôle} : le nom, le prix et la catégorie sont obligatoires ; le prix et la quantité doivent
        être positifs ou nuls ; la catégorie doit être connue. Chaque ligne invalide reçoit un motif précis
        (par exemple « prix négatif » ou « catégorie inconnue (jouets) »).
  \item \textbf{Chargement} : les lignes valides sont insérées dans \code{produits}, ou mises à jour si un
        produit de même nom et de même catégorie existe déjà, ce qui évite les doublons. Les lignes
        invalides sont écrites dans \code{produits_rejets}.
\end{enumerate}

\subsection{Le workflow d'orchestration}

Le workflow \code{wf_etl_produits} enchaîne la vérification des connexions, l'exécution du pipeline, le calcul
d'un bilan (pipeline \code{pl_bilan}) puis la journalisation. En cas d'erreur, il écrit un message d'échec et
se termine avec un code de sortie non nul : l'ETL peut ainsi être intégré à un test automatique.

\subsection{Exécution et résultats}

\begin{Code}
docker compose up -d            # démarre les bases
docker compose run --rm etl     # exécute le workflow puis s'arrête
\end{Code}

\begin{table}[H]
\centering\small
\begin{tabular}{|l|r|}
\hline
\rowcolor{bleu!15}
\textbf{Indicateur} & \textbf{Valeur} \\
\hline
Lignes lues dans MySQL & 40 \\
\hline
Lignes valides & 30 \\
\hline
Produits chargés dans PostgreSQL & 29 \\
\hline
Lignes rejetées & 10 \\
\hline
\end{tabular}
\caption{Résultats de l'ETL sur le jeu de test (source : rapport de la partie ETL, \code{etl/RAPPORT-ETL.md}).}
\end{table}

Les 10 lignes rejetées portent chacune un motif précis, relevé dans la table \code{produits_rejets} :

\begin{table}[H]
\centering\small
\begin{tabularx}{\textwidth}{|L{2.2cm}|Y|}
\hline
\rowcolor{bleu!15}
\textbf{Ligne source} & \textbf{Motif de rejet} \\
\hline
31, 32, 33 & nom manquant (valeur nulle, vide, ou espaces uniquement) \\
\hline
34 & prix manquant \\
\hline
35 & prix négatif \\
\hline
36 & quantité négative \\
\hline
37, 38 & catégorie manquante (valeur nulle ou vide) \\
\hline
39 & catégorie inconnue (jouets) \\
\hline
40 & nom manquant, prix manquant \\
\hline
\end{tabularx}
\caption{Motifs de rejet observés après exécution de l'ETL.}
\end{table}

L'écart entre 30 lignes valides et 29 produits chargés s'explique par le doublon : la ligne~30 met à jour
« Souris sans fil » au lieu de créer une nouvelle ligne. La catégorie « jouets », absente du référentiel,
est rejetée. Une seconde exécution donne les mêmes chiffres : l'ETL est \emph{idempotent}, car la table
des rejets est vidée à chaque passage et le chargement fonctionne par insertion ou mise à jour.

Les identifiants des bases ne sont écrits dans aucun fichier Hop : les connexions utilisent des variables
(\code{${MYSQL_PASSWORD}}, \code{${PG_PASSWORD}}\ldots) fournies par \code{.env} au démarrage du conteneur.

\capture{hop-pipeline.png}{Le pipeline \texttt{pl\_produits} dans l'interface Hop Web.}

% =====================================================================
\section{Partie 3 --- API REST Spring Boot}
% =====================================================================

\subsection{Technologies}

\begin{table}[H]
\centering\small
\begin{tabularx}{\textwidth}{|L{4.2cm}|Y|}
\hline
\rowcolor{bleu!15}
\textbf{Élément} & \textbf{Choix} \\
\hline
Langage et framework & Java 17, Spring Boot 3.3.5 \\
\hline
Accès aux données & Spring Data JPA, Hibernate, pilote PostgreSQL \\
\hline
Validation & Jakarta Bean Validation (\code{@Valid}, \code{@NotBlank}, \code{@PositiveOrZero}\ldots) \\
\hline
Documentation & springdoc-openapi 2.6.0 (contrat OpenAPI et Swagger UI) \\
\hline
Supervision & Spring Boot Actuator et Micrometer (registre Prometheus) \\
\hline
Confort de code & Lombok (accesseurs, constructeurs) \\
\hline
Tests & JUnit 5, Mockito, MockMvc, base H2 en mode PostgreSQL \\
\hline
\end{tabularx}
\caption{Pile technique de l'API.}
\end{table}

\subsection{Architecture en couches}

Le code suit une architecture en couches classique (figure~\ref{fig:couches}), chaque couche n'appelant
que la suivante.

\begin{figure}[H]
\centering
\begin{adjustbox}{max width=\textwidth}
\begin{tikzpicture}[node distance=0.55cm,boite/.append style={text width=2.3cm,minimum height=1.3cm}]
  \node[ext]  (cl) {Client HTTP};
  \node[boite,right=of cl] (co) {Controller\\ routes, statuts HTTP};
  \node[boite,right=of co] (se) {Service\\ règles métier, transactions};
  \node[boite,right=of se] (re) {Repository\\ Spring Data JPA};
  \node[ext,right=of re]   (db) {PostgreSQL};
  \draw[flux,<->] (cl)--(co); \draw[flux,<->] (co)--(se);
  \draw[flux,<->] (se)--(re); \draw[flux,<->] (re)--(db);
  \node[boite,below=0.9cm of se,text width=9.4cm,fill=rouge!6,draw=rouge!70]
        (ex) {\texttt{GlobalExceptionHandler} : toutes les erreurs (validation, introuvable, JSON mal formé,
              route inconnue) sont converties en un même format JSON};
  \draw[dashed,color=rouge!70,-{Stealth[length=2mm]}] (co.south) -- (co.south |- ex.north);
\end{tikzpicture}
\end{adjustbox}
\caption{Couches de l'API.}
\label{fig:couches}
\end{figure}

\begin{description}
  \item[\code{entity}] La classe \code{Produit} est liée à la table \code{produits}. Le champ
        \code{dateCreation} est posé par la base (\code{DEFAULT NOW()}) et relu après l'insertion ; l'API ne
        l'écrit jamais, ce qui garantit qu'un \code{PUT} ne peut pas la modifier.
  \item[\code{dto}] Des objets de transfert séparent le contrat de l'API du modèle de données :
        \code{ProduitRequest} (entrée, avec les règles de validation), \code{ProduitResponse} (sortie),
        \code{PageResponse} (résultat paginé) et \code{SuppressionResponse} (confirmation de suppression).
        Ce sont des \emph{records} Java, donc immuables.
  \item[\code{repository}] Interface Spring Data ; elle ajoute la recherche par catégorie sans tenir compte
        de la casse.
  \item[\code{service}] Contient la logique et les transactions : lecture seule par défaut, écriture
        explicite pour la création, la modification et la suppression.
  \item[\code{controller}] Déclare les routes sous \code{/api/v1/produits} et documente chaque réponse pour
        OpenAPI. La route \code{/} redirige vers Swagger UI.
  \item[\code{exception}] Un gestionnaire centralisé (\code{@RestControllerAdvice}) produit toutes les
        réponses d'erreur.
\end{description}

\subsection{Les endpoints}

\begin{table}[H]
\centering\small
\begin{tabularx}{\textwidth}{|L{2.5cm}|L{1.7cm}|L{4.2cm}|Y|}
\hline
\rowcolor{bleu!15}
\textbf{Opération} & \textbf{Méthode} & \textbf{URL} & \textbf{Réponses} \\
\hline
Lister       & GET    & \code{/api/v1/produits} & 200 ; 400 si \code{page} < 0 ou \code{size} hors de 1 à 100 \\
\hline
Consulter    & GET    & \code{/api/v1/produits/{id}} & 200 ; 404 \\
\hline
Créer        & POST   & \code{/api/v1/produits} & 201 avec l'en-tête \code{Location} ; 400 \\
\hline
Modifier     & PUT    & \code{/api/v1/produits/{id}} & 200 ; 400 ; 404 \\
\hline
Supprimer    & DELETE & \code{/api/v1/produits/{id}} & 200 avec message de confirmation ; 404 \\
\hline
\end{tabularx}
\caption{Ressource \code{produits}.}
\end{table}

La liste accepte les paramètres \code{categorie} (filtre, casse ignorée), \code{page} (défaut 0) et
\code{size} (défaut 20, maximum 100). Les résultats sont triés par identifiant et renvoyés sous la forme
\texttt{\{content, page, size, totalElements, totalPages\}}.

Trois endpoints techniques complètent l'API : \code{/actuator/health} (état de l'application),
\code{/actuator/prometheus} (métriques) et \code{/v3/api-docs} (contrat OpenAPI).

\subsection{Validation et gestion des erreurs}

Les règles de validation sont portées par \code{ProduitRequest} : le nom est obligatoire et limité à
150 caractères ; le prix est obligatoire, positif ou nul, avec au plus deux décimales ; la quantité est
obligatoire, positive ou nulle ; la catégorie est obligatoire et limitée à 100 caractères. Les paramètres
de pagination sont validés de la même façon.

Toutes les erreurs ont le \textbf{même format JSON}, y compris pour une route inconnue ou un corps JSON
mal formé, ce qui simplifie le travail des applications clientes :

\begin{Code}
{"timestamp":"2026-10-28T15:36:52Z","status":400,"error":"Bad Request",
 "message":"Données invalides","path":"/api/v1/produits",
 "details":[{"champ":"prix","message":"Le prix doit être positif ou nul"}]}
\end{Code}

\begin{table}[H]
\centering\small
\begin{tabularx}{\textwidth}{|L{5.6cm}|L{1.2cm}|Y|}
\hline
\rowcolor{bleu!15}
\textbf{Cause} & \textbf{Statut} & \textbf{Message} \\
\hline
Corps invalide (champ vide, prix négatif\ldots) & 400 & « Données invalides » + liste \code{details} \\
\hline
Paramètre de pagination hors limites            & 400 & « Paramètres invalides » + \code{details} \\
\hline
Identifiant non numérique (\code{/produits/abc}) & 400 & « Valeur invalide pour le paramètre 'id' » \\
\hline
Corps JSON absent ou mal formé                  & 400 & « Corps de requête JSON absent ou mal formé » \\
\hline
Produit inexistant                              & 404 & « Produit introuvable avec l'id \ldots » \\
\hline
Route inconnue                                  & 404 & « Ressource introuvable : \ldots » \\
\hline
\end{tabularx}
\caption{Erreurs gérées par \code{GlobalExceptionHandler}.}
\end{table}

\subsection{Contrat OpenAPI}

Le contrat est généré automatiquement par springdoc à partir du code puis \emph{exporté dans le dépôt}
(\code{api/openapi.yaml}) par un test dédié. Le contrat versionné permet de voir en une revue de code ce
qu'un changement modifie pour les clients. La documentation interactive est disponible sur
\code{/swagger-ui/index.html}.

\capture{swagger-ui.png}{Documentation interactive de l'API (Swagger UI).}

\subsection{Deux améliorations apportées aux réponses}

Après une première version, deux réponses jugées trop pauvres ont été enrichies.

\begin{itemize}
  \item \textbf{La suppression} renvoyait un \code{204} sans corps. Elle renvoie désormais un \code{200}
        avec un message de confirmation et le produit tel qu'il était avant suppression :
\begin{Code}
{"message":"Produit 11 supprimé",
 "produit":{"id":11,"nom":"Souris","prix":19.90,"quantite":40,
            "categorie":"Informatique","dateCreation":"2026-10-28T15:36:52"}}
\end{Code}
        Ce changement modifie le contrat (le statut passe de 204 à 200) ; les tests, le contrat OpenAPI,
        le README et la collection Bruno ont été mis à jour en conséquence.
  \item \textbf{L'état de santé} (\code{/actuator/health}) ne renvoyait que \texttt{\{"status":"UP"\}}. Avec
        \code{show-details: always}, il détaille chaque composant : base PostgreSQL, espace disque libre et
        \emph{ping}. Un incident de base de données devient visible sans consulter les journaux.
\end{itemize}

\subsection{Conteneurisation}

Le \code{Dockerfile} de l'API a deux étapes : une image Maven compile l'application, puis seule l'archive
\code{.jar} est copiée dans une image \emph{JRE 17} légère, exécutée par un utilisateur non privilégié. Les
dépendances sont téléchargées dans une couche séparée, ce qui accélère les reconstructions.

% =====================================================================
\section{Partie 4 --- Supervision avec Prometheus}
% =====================================================================

\subsection{Principe}

L'API expose ses métriques au format Prometheus sur \code{/actuator/prometheus}, grâce à Actuator et
Micrometer. Le serveur Prometheus les \emph{collecte} (\emph{scrape}) toutes les 5~secondes, les stocke
et permet de les interroger avec le langage PromQL, sur \code{http://localhost:9090}.

\begin{Code}
global:
  scrape_interval: 5s
scrape_configs:
  - job_name: "prometheus"
    static_configs:
      - targets: ["prometheus:9090"]
  - job_name: "produits-api"
    metrics_path: "/actuator/prometheus"
    static_configs:
      - targets: ["api:8080"]
\end{Code}

La cible de l'API est désignée par son nom de service Docker (\code{api:8080}) : la configuration est donc
identique sur toutes les machines. Dans l'interface de Prometheus (page \emph{Status~>~Targets}), les deux
cibles apparaissent à l'état \code{UP}.

\subsection{Métriques disponibles}

\begin{table}[H]
\centering\small
\begin{tabularx}{\textwidth}{|L{6.4cm}|Y|}
\hline
\rowcolor{bleu!15}
\textbf{Métrique} & \textbf{Information} \\
\hline
\code{http_server_requests_seconds_*} & Nombre et durée des requêtes, par méthode, URI, statut et résultat \\
\hline
\code{jvm_memory_used_bytes}          & Mémoire utilisée (tas et hors-tas) \\
\hline
\code{hikaricp_connections_*}         & Connexions actives et en attente vers PostgreSQL \\
\hline
\code{process_cpu_usage}              & Utilisation du processeur \\
\hline
\code{up}                             & 1 si la cible répond à la collecte, 0 sinon \\
\hline
\end{tabularx}
\caption{Principales familles de métriques.}
\end{table}

\subsection{Requêtes PromQL utiles}

\begin{Code}
up                                                    # état de chaque cible
sum by (uri) (rate(http_server_requests_seconds_count{uri=~"/api/.*"}[5m]))
                                                      # requêtes par seconde et par route
sum by (status) (http_server_requests_seconds_count{uri=~"/api/.*"})
                                                      # nombre de réponses par code HTTP
sum by (area) (jvm_memory_used_bytes)                 # mémoire de la JVM
hikaricp_connections_active                           # connexions actives à la base
\end{Code}

Ces requêtes ont été exécutées sur l'environnement du projet et renvoient des séries de données.
Elles permettent par exemple de repérer une hausse des erreurs 4xx/5xx ou une saturation de la mémoire.

\capture{prometheus-targets.png}{Cibles de collecte dans Prometheus.}
\capture{prometheus-requetes.png}{Requête PromQL sur le trafic de l'API.}

% =====================================================================
\section{Partie 5 --- Tests et validation}
% =====================================================================

La qualité est vérifiée à quatre niveaux, du plus fin au plus global : tests Java (unitaires et d'intégration), collections Postman et Bruno sur l'API en marche, puis script de bout en bout sur toute la chaîne.

\subsection{Tests unitaires du service}

Neuf tests JUnit~5 isolent la couche service en simulant le repository avec Mockito. Ils vérifient la
création (avec normalisation des espaces), la modification, la suppression (y compris la valeur
renvoyée), la liste avec ou sans filtre de catégorie, et le déclenchement de
\code{ProduitNotFoundException} quand un produit est absent.

\subsection{Tests d'intégration de l'API}

Dix-huit tests parcourent l'application de bout en bout --- contrôleur, validation, service, JPA --- avec
MockMvc et une base H2 configurée en mode PostgreSQL. Le schéma de test
(\code{src/test/resources/schema-test.sql}) reproduit celui de PostgreSQL, y compris les contraintes
\code{CHECK}. Les cas couverts sont notamment :
\begin{itemize}
  \item création valide (201, en-tête \code{Location}, date de création posée) ;
  \item création invalide (400, champs signalés), prix manquant et JSON mal formé ;
  \item modification, dont la conservation de la date de création, et modification invalide ;
  \item suppression (200 avec confirmation), puis disparition du produit (404) ;
  \item consultation d'un produit absent (404) et identifiant non numérique (400) ;
  \item pagination, filtre de catégorie, paramètres de pagination hors limites ;
  \item route inconnue (404 au format JSON), redirection de la racine vers Swagger, exposition du contrat
        OpenAPI et de l'état de santé.
\end{itemize}
Ces 27 tests (9 + 18) s'exécutent sans Docker avec \code{mvn test} ; ils réussissent tous. Un 28\textsuperscript{e}
test, qui exporte le contrat OpenAPI, est désactivé par défaut et se lance avec
\code{-Dopenapi.export=true}. Le projet cible Java~17 : avec une version plus récente (par exemple Java~27),
la compilation échoue à cause de Lombok.

\subsection{Collection Bruno}

La collection du dossier \code{bruno/} teste l'API réelle, en marche, par HTTP. Elle contient
\textbf{18 requêtes} en trois dossiers, chacune avec des assertions et des tests aux messages explicites.

\begin{table}[H]
\centering\small
\begin{tabularx}{\textwidth}{|L{2.6cm}|L{1.9cm}|Y|}
\hline
\rowcolor{bleu!15}
\textbf{Dossier} & \textbf{Requêtes} & \textbf{Contenu} \\
\hline
Produits    & 7 & Lister, lister par catégorie, créer, consulter, modifier, supprimer, vérifier la suppression \\
\hline
Erreurs     & 8 & Données invalides (création et modification), pagination invalide, identifiant non numérique, produit inexistant (consulter, modifier, supprimer), route inconnue \\
\hline
Supervision & 3 & Santé détaillée, métriques Prometheus, contrat OpenAPI \\
\hline
\end{tabularx}
\caption{Organisation de la collection Bruno.}
\end{table}

Le dossier \emph{Produits} forme un scénario : la requête de création mémorise l'identifiant du produit dans
une variable, que les requêtes suivantes réutilisent. Le parcours complet crée puis supprime un produit,
et laisse donc la base dans son état initial. La collection s'ouvre dans Bruno par \emph{Open Collection},
avec l'environnement \emph{Local} (\code{baseUrl = http://localhost:8080}).

\capture{bruno-runner.png}{Exécution de la collection Bruno avec le Runner.}

\subsection{Collection Postman}

L'énoncé demande de tester l'API avec Postman. Le dossier \code{postman/} contient donc la même collection
sous forme de fichier Postman (\code{API-Produits.postman_collection.json}) : mêmes 18 requêtes, mêmes trois
dossiers, mêmes vérifications. Elle s'importe dans Postman (\emph{Import}) et se lance avec le
\emph{Collection Runner}, ou en ligne de commande avec \code{newman}. La variable \code{baseUrl} vaut
\code{http://localhost:8080} et la requête de création mémorise l'identifiant du produit, comme dans Bruno.
Exécutée avec \code{newman} sur l'API en marche le 29 octobre 2026 : \textbf{18 requêtes, 41 assertions, 0 échec}.

\subsection{Test de bout en bout}

Le test de bout en bout vérifie que toute la chaîne fonctionne ensemble, depuis une machine vierge.
Le scénario est le suivant ; les valeurs attendues proviennent du jeu de test (section~\ref{sec:jeu-ref}).

\begin{table}[H]
\centering\small
\begin{tabularx}{\textwidth}{|L{0.8cm}|L{7.4cm}|Y|}
\hline
\rowcolor{bleu!15}
\textbf{N°} & \textbf{Action} & \textbf{Résultat attendu} \\
\hline
1 & \code{docker compose down -v} puis \code{docker compose up -d --build} & 4 services démarrés ; MySQL et PostgreSQL \code{healthy} \\
\hline
2 & \code{bash db/verifier-bases.sh} & 40 lignes dans MySQL ; \code{produits} et \code{produits_rejets} vides \\
\hline
3 & \code{docker compose run --rm etl} & Code de sortie 0 ; 40 lues, 29 chargées, 10 rejetées \\
\hline
4 & \code{GET /api/v1/produits?size=100} & \code{totalElements = 29} \\
\hline
5 & \code{GET /api/v1/produits?categorie=informatique} & Uniquement des produits « Informatique », quelle que soit la casse demandée \\
\hline
6 & Dossier \emph{Produits} de Bruno & Création, lecture, modification, suppression, puis 404 \\
\hline
7 & Dossier \emph{Erreurs} de Bruno & Tous les cas d'erreur renvoient le format JSON standard \\
\hline
8 & Page \emph{Targets} de Prometheus & Cible \code{produits-api} à l'état \code{UP} ; compteurs des requêtes en hausse \\
\hline
9 & Second passage de l'ETL & Mêmes chiffres (29 produits), sans doublon \\
\hline
\end{tabularx}
\caption{Scénario de bout en bout et résultats attendus.}
\end{table}

\paragraph{Résultats observés.} Le scénario a été exécuté le 28 octobre 2026 sur une machine neuve côté
Docker (volumes supprimés par \code{docker compose down -v}). Les 9 étapes réussissent (tableau~\ref{tab:resultats}).

\begin{table}[H]
\centering\small
\begin{tabularx}{\textwidth}{|L{0.8cm}|Y|L{1.5cm}|}
\hline
\rowcolor{bleu!15}
\textbf{N°} & \textbf{Résultat observé} & \textbf{Statut} \\
\hline
1 & Les 4 services permanents démarrent en une vingtaine de secondes ; MySQL et PostgreSQL sont \code{healthy} après environ 9\,s ; l'API répond \code{UP}. & \textbf{OK} \\
\hline
2 & MySQL contient 40 lignes ; \code{produits} et \code{produits_rejets} sont vides ; l'API renvoie \code{totalElements = 0}. & \textbf{OK} \\
\hline
3 & Code de sortie 0. Journal : 40 lignes lues, 29 produits chargés, 10 rejetées. Le workflow dure 1,6\,s ; la commande complète 1\,min\,23\,s à la première exécution (construction de l'image Hop comprise). & \textbf{OK} \\
\hline
4 & \code{totalElements = 29}, produits triés par identifiant croissant. & \textbf{OK} \\
\hline
5 & Le filtre renvoie 6 produits « Informatique » pour \code{informatique}, \code{INFORMATIQUE} et \code{Informatique}. Répartition des 29 produits : Alimentation~5, Electronique~7, Informatique~6, Maison~5, Vetements~6 ; aucun produit « jouets ». & \textbf{OK} \\
\hline
6 & Parcours création, lecture, modification, suppression : 201 avec \code{Location}, 200, 200 (\code{dateCreation} inchangée), 200 avec message « Produit 31 supprimé », puis 404 ; la base retrouve ses 29 produits. 9 vérifications. & \textbf{OK} \\
\hline
7 & Les 8 cas d'erreur renvoient le statut attendu (400 ou 404) et le format JSON standard. 8 vérifications. & \textbf{OK} \\
\hline
8 & Cibles \code{produits-api} et \code{prometheus} à l'état \code{UP} ; le compteur des requêtes sur \code{/api} passe de 29 à 31 après deux requêtes supplémentaires. & \textbf{OK} \\
\hline
9 & Code de sortie 0 ; mêmes chiffres (40, 29, 10) ; l'API liste toujours 29 produits, soit 29 couples (nom, catégorie) distincts : aucun doublon. Commande complète : 18\,s. & \textbf{OK} \\
\hline
\end{tabularx}
\caption{Résultats du scénario de bout en bout (28 vérifications automatisées sur les étapes 4 à 7 et la supervision, 0 échec).}
\label{tab:resultats}
\end{table}

Les étapes 4 à 7 ont été rejouées avec la collection Bruno exécutée en ligne de commande
(\code{bru run --env Local}, dossiers \emph{Produits}, \emph{Erreurs} et \emph{Supervision}) :
\textbf{18 requêtes réussies, 33 tests et 19 assertions sans échec}.

\subsection{Script automatisé de bout en bout}

Pour que ce scénario soit rejouable par n'importe quel relecteur, il a été automatisé dans le script
\code{tests/e2e.sh}. Il suffit de lancer \code{bash tests/e2e.sh --clean} : le script réinitialise les
volumes Docker, démarre les services, exécute l'ETL, interroge l'API et Prometheus, puis affiche
\texttt{[OK]} ou \texttt{[KO]} pour chacun des 32 contrôles. Son code de retour vaut 0 si tout réussit, ce
qui le rend utilisable dans une chaîne d'intégration continue. Contrairement aux collections Postman et
Bruno, qui ne testent que l'API, il vérifie aussi directement le contenu des deux bases de données.

\begin{table}[H]
\centering\small
\begin{tabularx}{\textwidth}{|L{3.6cm}|Y|L{2.4cm}|}
\hline
\rowcolor{bleu!15}
\textbf{Étape} & \textbf{Contrôles} & \textbf{Résultat} \\
\hline
Démarrage et source & API \code{UP} ; 40 lignes dans \code{produits_source} & 2 OK \\
\hline
ETL & Code de sortie 0 ; 29 produits chargés ; 10 rejets, tous avec un motif ; aucune catégorie hors référentiel ; aucun prix négatif ; relance sans doublon (toujours 29) & 7 OK \\
\hline
Lecture par l'API & Liste (\code{totalElements = 29}) ; filtre par catégorie égal au résultat de la requête SQL & 4 OK \\
\hline
Création, modification, suppression & 201 ; produit présent dans PostgreSQL ; 200 ; nom modifié en base ; 200 ; 404 après suppression ; produit absent de PostgreSQL ; retour à 29 produits & 10 OK \\
\hline
Erreurs & Corps invalide (400, message standard) ; identifiant inexistant (404) ; identifiant non numérique (400) ; pagination invalide (400) & 5 OK \\
\hline
Supervision & \code{/actuator/prometheus} (200, métrique HTTP présente) ; cible \code{produits-api} à \code{up = 1} ; requêtes HTTP comptées par Prometheus & 4 OK \\
\hline
\end{tabularx}
\caption{Contrôles du script \code{tests/e2e.sh} (le détail est dans \code{tests/RESULTATS-ATTENDUS.md}).}
\end{table}

\paragraph{Résultat observé.} Exécuté le 29 octobre 2026 avec l'option \code{--clean} (volumes supprimés,
images reconstruites) : \textbf{32 contrôles réussis, 0 échec}.

% =====================================================================
\section{Difficultés rencontrées et solutions}
\label{sec:difficultes}
% =====================================================================

\begin{table}[H]
\centering\small
\begin{tabularx}{\textwidth}{|L{4.2cm}|Y|}
\hline
\rowcolor{bleu!15}
\textbf{Difficulté} & \textbf{Solution} \\
\hline
Talend Open Studio n'est plus maintenu (2024) &
  Choix d'Apache Hop, documenté et justifié, conforme à l'énoncé qui autorise une alternative \\
\hline
Apache Hop ne fournit pas le pilote MySQL (licence) &
  Pilote ajouté automatiquement à la construction de l'image Docker \\
\hline
Aucun composant Hop ne supprime les accents &
  Court script JavaScript intégré au pipeline \\
\hline
L'image officielle de Hop affichait les mots de passe dans son journal &
  Identifiants transmis par un fichier temporaire créé au démarrage du conteneur \\
\hline
Ports déjà occupés sur certaines machines (PostgreSQL local sur 5432, autre conteneur sur 5433) &
  Ports publiés configurables dans \code{.env} (\code{PG_PORT}) ; l'API, dans Docker, joint toujours \code{postgres:5432} en interne \\
\hline
Compilation Maven en échec avec Java~27 (Lombok incompatible) &
  Utilisation du JDK~17 visé par le projet ; l'image Docker n'est pas concernée \\
\hline
Réponses de l'API jugées trop pauvres (suppression sans corps, santé réduite à \texttt{UP}) &
  Suppression avec confirmation et santé détaillée, tests et contrat mis à jour \\
\hline
\end{tabularx}
\caption{Principales difficultés et réponses apportées.}
\end{table}

% =====================================================================
\section{Organisation du travail}
% =====================================================================

Le travail a été réparti en cinq parties, une par membre. Le dépôt Git applique un flux simple : une
branche par fonctionnalité (\code{feat/*}), intégrée dans \code{main} par une demande de fusion
(\emph{pull request}).

\begin{table}[H]
\centering\small
\begin{tabularx}{\textwidth}{|L{1.1cm}|L{2.6cm}|Y|L{5.0cm}|}
\hline
\rowcolor{bleu!15}
\textbf{Partie} & \textbf{Membre} & \textbf{Réalisations} & \textbf{Branche Git} \\
\hline
1 & BOUGOUMA Leïla & Scripts SQL des deux bases, jeu de test de 40 lignes (\code{DONNEES-TEST.md}), \code{docker-compose.yml}, \code{.env.example}, script \code{verifier-bases.sh} & \code{main} (commit initial) \\
\hline
2 & OUARMA Landry & Projet Apache Hop (pipelines \code{pl_produits} et \code{pl_bilan}, workflow \code{wf_etl_produits}), référentiel de catégories, services Docker \code{etl} et \code{hop-web}, installation et rapport de l'ETL & \code{feat/etl-hop} (PR~\#3) \\
\hline
3 & WAONGO Osiris & API Spring Boot en couches, validation, gestion centralisée des erreurs, contrat OpenAPI, \code{Dockerfile}, tests unitaires et d'intégration & \code{feat/api-dev} (PR~\#2) \\
\hline
4 & ZAPRE Malika & Fichier \code{prometheus.yml}, service Prometheus dans Docker Compose, exposition des métriques de l'API & \code{feat/monitoring-prometheus} (PR~\#1) \\
\hline
5 & KOURA Stéphane & Collection Bruno (18 requêtes), enrichissement des réponses de l'API (suppression, santé détaillée), scénario de bout en bout, rapport LaTeX et guide de présentation & \code{feat/tests} (PR~\#4) \\
\hline
\end{tabularx}
\caption{Répartition des parties.}
\end{table}

% =====================================================================
\section{Limites et perspectives}
% =====================================================================

\begin{itemize}
  \item \textbf{Sécurité} : l'API n'a pas d'authentification, et l'état de santé détaillé est public. C'est
        acceptable pour un environnement de développement ; en production, il faudrait ajouter Spring
        Security (jetons JWT) et limiter l'accès aux endpoints \code{/actuator}.
  \item \textbf{ETL automatique} : l'ETL se lance à la demande. Il pourrait être planifié (tâche
        périodique) ou déclenché à l'arrivée de nouvelles données.
  \item \textbf{Supervision} : ajouter des règles d'alerte (taux d'erreurs, indisponibilité) avec
        Alertmanager, et des tableaux de bord Grafana.
  \item \textbf{Tests} : utiliser Testcontainers pour tester l'API contre un vrai PostgreSQL plutôt que H2,
        et exécuter le script \code{tests/e2e.sh} dans une chaîne d'intégration continue à chaque modification.
  \item \textbf{Évolution de l'API} : versionnement déjà prévu par le préfixe \code{/api/v1} ; ajout possible
        de la recherche par nom, du tri configurable et de la suppression logique.
\end{itemize}

% =====================================================================
\section{Conclusion}
% =====================================================================

Le projet atteint son objectif : une chaîne complète, reproductible et vérifiée, qui transforme des données
brutes et imparfaites en un service web fiable et supervisé. L'ETL garantit la qualité des données, la base
cible la protège par des contraintes, l'API l'expose avec un contrat clair et des erreurs homogènes, et
Prometheus permet de suivre son comportement. La combinaison de Docker Compose, d'un contrat OpenAPI
versionné, de 27 tests automatisés, de collections Postman et Bruno et d'un script de bout en bout rend le projet facile à lancer, à relire et à
faire évoluer. Le travail en groupe, avec des parties indépendantes reliées par des conventions communes
(schéma SQL, variables d'environnement, ports), a permis d'avancer en parallèle sans blocage.

\appendix
% =====================================================================
\section{Lancer le projet}
% =====================================================================

\begin{Code}
# 1. Configuration (changer PG_PORT si le port 5432 est déjà pris)
cp .env.example .env

# 2. Démarrer MySQL, PostgreSQL, l'API et Prometheus
docker compose up -d --build
docker compose ps                      # attendre "healthy"

# 3. Charger les données propres dans PostgreSQL
docker compose run --rm etl

# 4. Vérifier
curl http://localhost:8080/actuator/health
curl "http://localhost:8080/api/v1/produits?size=100"

# 5. Test de bout en bout (32 contrôles attendus)
bash tests/e2e.sh
\end{Code}

\begin{table}[H]
\centering\small
\begin{tabularx}{\textwidth}{|L{8.9cm}|Y|}
\hline
\rowcolor{bleu!15}
\textbf{Adresse} & \textbf{Contenu} \\
\hline
\code{http://localhost:8080/swagger-ui/index.html} & Documentation interactive de l'API \\
\hline
\code{http://localhost:8080/actuator/health}      & État de l'application et de ses composants \\
\hline
\code{http://localhost:9090}                      & Prometheus \\
\hline
\code{http://localhost:8090/ui}                   & Hop Web (profil \code{hop-web}) \\
\hline
\end{tabularx}
\caption{Adresses utiles.}
\end{table}

% =====================================================================
\section{Structure du dépôt}
% =====================================================================

\begin{Code}
.
|-- docker-compose.yml      services : mysql, postgres, etl, hop-web, api, prometheus
|-- .env.example            modèle de configuration (le .env n'est pas versionné)
|-- db/                     scripts SQL des deux bases, jeu de test, verifier-bases.sh
|-- etl/                    projet Apache Hop, référentiel de catégories, Dockerfile
|-- api/                    API Spring Boot : src/, pom.xml, Dockerfile, openapi.yaml
|-- monitoring/             prometheus.yml
|-- postman/                collection Postman (18 requêtes)
|-- bruno/                  collection Bruno (18 requêtes)
|-- tests/                  e2e.sh (bout en bout) et résultats attendus
|-- presentation/          présentation PowerPoint
|-- rapport/                ce rapport (LaTeX)
`-- docs/                   guide de préparation de la présentation
\end{Code}

\section{Exemples de requêtes et de réponses}

\paragraph{Créer un produit.}
\begin{Code}
curl -X POST http://localhost:8080/api/v1/produits \
  -H 'Content-Type: application/json' \
  -d '{"nom":"Clavier","description":"AZERTY","prix":59.90,
       "quantite":25,"categorie":"Informatique"}'

HTTP/1.1 201
Location: http://localhost:8080/api/v1/produits/12
{"id":12,"nom":"Clavier","description":"AZERTY","prix":59.90,"quantite":25,
 "categorie":"Informatique","dateCreation":"2026-10-28T15:36:52.77"}
\end{Code}

\paragraph{Lister avec filtre et pagination.}
\begin{Code}
GET /api/v1/produits?categorie=informatique&page=0&size=2

{"content":[ {...}, {...} ],"page":0,"size":2,"totalElements":7,"totalPages":4}
\end{Code}

\paragraph{Produit inexistant.}
\begin{Code}
GET /api/v1/produits/999999            ->  404

{"timestamp":"2026-10-28T15:36:52Z","status":404,"error":"Not Found",
 "message":"Produit introuvable avec l'id 999999",
 "path":"/api/v1/produits/999999"}
\end{Code}

\end{document}
